\documentclass[10pt,twocolumn]{article}
\usepackage[a4paper,margin=0.8in,columnsep=0.28in]{geometry}
\usepackage[T1]{fontenc}
\usepackage{lmodern}
\usepackage{microtype}
\usepackage{booktabs}
\usepackage{amsmath}
\usepackage{xcolor}
\usepackage[hidelinks]{hyperref}
\usepackage{titlesec}
\titleformat{\section}{\large\bfseries}{\thesection}{0.6em}{}
\titleformat{\subsection}{\normalsize\bfseries}{\thesubsection}{0.6em}{}
\newcommand{\us}{\,\textmu s}
\newcommand{\ms}{\,ms}
\newcommand{\pol}[1]{\texttt{#1}}
\newcommand{\finding}[1]{\textbf{#1}}

\title{\textbf{What Should the Scheduler Know?}\\[3pt]
\large Compiler-Derived Behavioral Hints in a From-Scratch x86-64 Operating System}
\author{Teye John Harrison Junior\thanks{AI tools assisted with parts of the implementation,
the experiment tooling and the drafting of this report; the research questions,
design decisions and interpretation are the author's. All source code, raw data, analysis
scripts and pre-registered predictions: \url{https://github.com/JohnTeye0000/aetheros}.}}
\date{Technical report, \today}

\begin{document}
\twocolumn[
\maketitle
\begin{abstract}
\noindent
A CPU scheduler must guess how each program will behave, so it infers that behavior from
how the program has run so far. A compiler sees the program \emph{before} it runs. This work
asks whether static knowledge from a compiler, delivered to the kernel as a declared
behavioral hint, can replace or improve runtime inference, and when it should not be trusted.
I built the complete stack from scratch: a preemptive x86-64 kernel in Rust; a native code
generator for PlainLang, an English-like language, differentially tested on 232 programs
inside the kernel; a static classifier that labels programs CPU-bound, I/O-bound or
interactive; and a deterministic benchmark harness under instruction-counted emulation.
Across more than 3{,}400 runs and six policies, hint-driven scheduling reduced interactive
wake-up latency from 7--10\ms{} (round-robin, MLFQ) to about 2\us. However, a CFS-style fair
scheduler with \emph{no} hints matched it. Hints helped measurably with efficiency, cutting
context switches by up to $14\times$ at unchanged completion time. Used as priorities, wrong
hints were costly (106\ms); a hybrid that uses hints only to size time slices held about
3\us{} even with every hint wrong. The classifier matched measured behavior on 76\% of 150
random programs, and every error underestimated demand. In the kernel, the direction of an
error determined who paid for it. Predictions were committed before each experiment, and
failures are reported alongside successes.
\end{abstract}
\vspace{1.2em}
]

\section{Introduction}

Operating-system schedulers decide, many times per second, which program runs next.
They usually cannot know what a program will do, so they infer it from history.
Multi-level feedback queues (MLFQ) demote threads that use their whole time slice and
favor those that block early~\cite{corbato62,ostep}. Fair schedulers such as Linux's
Completely Fair Scheduler (CFS) track how much CPU time each thread has received and run
the least-served one~\cite{cfs}. Both learn from behavior after the fact.

A compiler has a different vantage point. It sees loops, their bounds, and the places where
a program waits, before the program has executed a single instruction. That suggests a
question that is easy to state and harder to answer honestly:

\begin{quote}
\emph{Can compiler-derived information about a program's behavior make CPU scheduling
better, and under what conditions does it make it worse?}
\end{quote}

\noindent Answering this needs control over every layer at once: the language, the
compiler, the program format that carries the hint, the kernel that schedules, and the
instrumentation that measures. No off-the-shelf stack gives that, so I built one.

\paragraph{Contributions.}
\begin{enumerate}\setlength\itemsep{1pt}
\item A full experimental platform: a from-scratch x86-64 kernel with six switchable
scheduling policies and microsecond-resolution per-thread accounting, a compiler from
PlainLang to native x86-64, a static behavior classifier, and an unattended,
deterministic benchmark harness (\S\ref{sec:system}).
\item Evidence that hints beat textbook schedulers by three to four orders of magnitude on
interactive latency, but that a hint-free fair scheduler matches them; the measurable value
of hints is fewer context switches (\S\ref{sec:results}).
\item A hybrid design, \pol{hfair}, in which hints influence only time-slice length. It
keeps fair scheduling's latency, recovers most of the switch savings, and cannot be harmed
by a wrong hint.
\item An accuracy study of the compiler's classifier on 150 random programs, showing a
one-directional error bias, and a kernel experiment showing who bears the cost of each
kind of error (\S\ref{sec:compiler}).
\item A pre-registered methodology in which predictions were written and committed before
each experiment, with every failed prediction reported (\S\ref{sec:method}).
\end{enumerate}

\section{System}\label{sec:system}

\subsection{The kernel}
AetherOS is written in Rust without the standard library and boots on x86-64. It provides
per-process page tables, ring-3 user programs, system calls through both \texttt{int 0x80}
and \texttt{SYSCALL/SYSRET}, NX, SMEP and SMAP protection, a small on-disk file system, and
preemptive scheduling driven by a 1\,kHz timer. Each context switch saves and restores the
SSE register file, verified with a negative control. Blocking uses real wait queues rather
than polling, and an idle thread never consumes a time slice. The scheduler records, for
every thread, its response time (creation to first dispatch), turnaround, CPU time, time spent
ready but not running, and the latency of every single wake-up (woken to actually running).

\subsection{Scheduling policies}
All six policies are switchable at run time (Table~\ref{tab:policies}). The base quantum is
$Q=10$\ms{} unless stated.

\begin{table}[t]
\centering\small
\begin{tabular}{@{}lp{5.2cm}@{}}
\toprule
Policy & Decision rule \\ \midrule
\pol{rr} & One queue, fixed quantum $Q$. \\
\pol{mlfq} & 3 levels, slice $Q\cdot2^{\ell}$; demote after using a level's allotment;
periodic boost to the top (200\ms). \\
\pol{hint} & Priority from the declared hint (interactive $>$ io $>$ cpu); threads waiting
$>100$\ms{} are aged to the top. \\
\pol{hmlfq} & MLFQ whose \emph{starting} level comes from the hint. \\
\pol{fair} & CFS-style: run the thread with least virtual runtime; slice
$\max(Q/n,1\,\text{ms})$; sleeper credit; wake-up preemption. No hints. \\
\pol{hfair} & \pol{fair}, but a thread hinted \texttt{cpu} gets a slice $k\times$ longer
(default $k=4$). Hints never affect \emph{who} runs. \\ \bottomrule
\end{tabular}
\caption{Scheduling policies.}\label{tab:policies}
\end{table}

\subsection{From source code to hint}
PlainLang is an English-like teaching language (\texttt{repeat 30 times}, \texttt{wait
0.004}). I wrote a backend that emits x86-64 machine code directly into the kernel's program
format, with the hint embedded in the file header. Its correctness was established by
differential testing: 232 programs, including 200 randomly generated ones, produce identical
output inside the kernel and in the reference interpreter.

The classifier estimates a program's \emph{work} (operations, multiplied through loop trip
counts) and its number of \emph{waits}, and labels it by work per wait: interactive below
2{,}000 operations, I/O-bound below 100{,}000, CPU-bound above. Its \emph{counted} analysis
adds constant propagation and a conservative recognizer for counted \texttt{while} loops.
Loops whose length it cannot determine are assumed to run 100 times. Every hint comes with a
written rationale, so a wrong guess can be traced to the rule that produced it.

\subsection{Measurement}
A host-side harness writes a plan of benchmark lines onto a virtual disk, boots the kernel
unattended in QEMU, and parses its serial output. QEMU runs with
\texttt{-icount shift=3,sleep=off}, so guest time advances 8\,ns per executed instruction,
independent of host load. Two boots with identical inputs agreed within 0.1\%. Workloads are
built from three synthetic programs: \texttt{wcpu} (pure computation), \texttt{wio} (30
rounds of $\sim$4.8\ms{} work then a 4\ms{} wait) and \texttt{wint} (30 rounds of a 15\ms{}
wait then $\sim$0.5\ms{} work), plus PlainLang programs whose hints come from the compiler.
Statistics are means with 95\% confidence intervals over seeds (10 for main experiments,
3--5 for sweeps). The primary dataset uses the optimized kernel build. An earlier
unoptimized build overstated scheduling costs by $10\times$ and is not used for any claim.

\section{Methodology}\label{sec:method}
For every experiment after the first round, predictions were written into the plan file's
header and committed to version control \emph{before} the experiment ran. They were then
scored by scripts rather than by eye. Where an explanation was developed after seeing a
failure, it was treated as a hypothesis and given its own pre-registered test. Each
measurement problem found along the way is documented, and affected data was discarded and
rerun. These included a starved benchmark launcher, an $O(n)$ frame allocator that made the
first run of each boot slow, a parser that dropped one row in 24, and a boost timer that
carried state between runs (\S\ref{sec:threats}).

\section{Results}\label{sec:results}

\subsection{Against textbook schedulers}
\finding{Correct hints cut interactive latency by three to four orders of magnitude.}
On a workload of two CPU-bound processes, one I/O-style and one interactive, the
interactive process's p95 wake-up latency is 2\us{} under \pol{hint}, against 9.8\ms{}
under \pol{rr} and 6.8\ms{} under \pol{mlfq}. A woken interactive thread has strictly
higher priority and preempts within one timer tick.

\finding{The MLFQ result depends on a constant I chose.} Sweeping the boost period showed
that MLFQ's periodic boost \emph{hurts} interactive latency, the opposite of my prediction.
The boost returns every CPU-bound thread to the interactive thread's level. Without the
boost, MLFQ reaches 1--2\us{}, but the I/O-style process sinks to 76\ms{} p95. A
comparison against a single MLFQ configuration would have been misleading.

\finding{Wrong hints are worse than none.} With deliberately inverted hints, \pol{hint}
gives the interactive process 106\ms{} p95 latency, bounded only by the 100\ms{} aging
threshold. Across threshold settings, latency tracked the threshold (27, 67, 106,
206\ms). With aging disabled, the victim was \emph{starved}: its first dispatch came after
1.29\,s. Its wake-up latency still looked small (p95 0.8\ms), because wake-up latency cannot
see a thread that never ran. All reports now include response time and a starvation audit.

\subsection{Against a fair scheduler}
\finding{A hint-free fair scheduler matches hints on interactive latency.} \pol{fair}
achieves 3\us{} against \pol{hint}'s 2\us. With many competitors, \pol{fair} is
better at the default aging setting (5\us{} vs.\ 3.6\ms{} with ten CPU-bound
processes). A sweep showed the hint scheduler's degradation is caused entirely by its aging
rule, which lifts a waiting CPU-bound thread straight to the interactive level. Replacing
it with graded aging (one level per 100\ms{} waited) removed that tail (3.6\ms{} to
4\us) and two other anomalies traced to the same rule. The cost: a victim of wrong hints
waits about twice as long before rescue (206\ms{} instead of 106\ms), though it is never
starved. The mechanism behind \pol{fair}'s result is simple: a thread that sleeps has used little CPU, so its virtual
runtime is small and it is picked first. The scheduler measures what the hint declares.
\pol{fair} also handles programs that break the hint scheduler. For a program interactive in
one phase and CPU-bound in the next, \pol{hint} gives 65\ms{} latency and \pol{fair}
0.002\ms.

\finding{What hints still buy is fewer context switches.} Makespan is within about 1\% for
every policy, so total throughput does not change. Switch counts differ greatly
(Table~\ref{tab:switches}). Because the emulator has no caches or TLB, real switches cost
more than measured here, which makes this the result most worth carrying to real hardware.
Lower mean turnaround for CPU-bound work under \pol{hint} is an ordering effect, not a
throughput gain. One qualification matters: \pol{fair}'s switch count is itself set by its
minimum slice. Raising that floor from 1 to 5\ms{} cut switches to $0.21\times$ with
interactive latency unchanged (3--6\us). Hints therefore buy switch savings \emph{without}
lengthening every slice, not savings that are unavailable any other way.
The same holds inside \pol{hint}: with equal slices at every level, its switch count on
four CPU-bound programs rose from 85 to 267, matching round robin (268). The entire
switch advantage comes from \emph{which} threads receive long slices, not from priority.
That is precisely the part \pol{hfair} retains.

\begin{table}[t]
\centering\small
\begin{tabular}{@{}lrrrr@{}}
\toprule
Workload & \pol{rr} & \pol{fair} & \pol{hint} & \pol{hfair}$_{16}$ \\ \midrule
4 CPU-bound & 268 & 1{,}295 & 86 & 92 \\
mixed (\texttt{all}) & 201 & 419 & 153 & 157 \\
6 CPU + 1 interactive & 437 & 3{,}938 & 248 & 293 \\ \bottomrule
\end{tabular}
\caption{Context switches per run (means).}\label{tab:switches}
\end{table}

\subsection{A hybrid: hints size slices, measurement picks who runs}
In \pol{hfair}, a \texttt{cpu} hint only lengthens that thread's slice. Selection is still
by virtual runtime, so a hint cannot delay anyone else by more than the slice factor. With
multiplier 16, switch counts come within 3--18\% of \pol{hint}'s while interactive latency
stays at 3--4\us. With \emph{wrong} hints, interactive latency stays at 3\us{} (\pol{hint}:
106\ms) and switch counts fall back to \pol{fair}'s level. A wrong hint removes the
optimization and nothing else. On PlainLang programs with compiler-derived hints,
\pol{hfair} gives 0.003\ms{} interactive latency with about 45\% fewer switches than
\pol{fair}.

\subsection{Out of sample: random workloads}
Because \pol{hfair} was designed after seeing the data it explains, I tested it on 30
randomly drawn workloads (1--8 CPU-bound, 0--3 I/O-style, 0--3 interactive processes; 720
runs per batch) with predictions committed first. The structural claims replicated: the
ordering $\text{switches}(\pol{hint}) \le \text{switches}(\pol{hfair}) \le
1.05\cdot\text{switches}(\pol{fair})$ held in all 30 mixes in both batches, and wrong hints
broke \pol{hint} in every mix with both CPU-bound and interactive processes (up to 164\ms{}
response time). Makespan stayed within 1.14\%. Three latency predictions failed. Every
failure involved several interactive processes, and the explanation I proposed afterwards
was given its own test. Wake-ups that fall in the same 1\ms{} timer tick must queue behind
each other whatever the policy. With no competing processes at all, two interactive
programs had 3.3\% of wake-ups delayed and three had 5.6--5.8\%, identically under all five
policies. Turning off aging restored \pol{hint}'s latency to 3--4\us{} in every mix with at
most two interactive processes. All four predictions of that follow-up held. The lesson:
p95 latency is a fair summary only with one interactive process, because with several there
is a floor that belongs to the workload, not to any scheduler.

\section{How good is the compiler?}\label{sec:compiler}

\subsection{Accuracy on random programs}
Five hand-written programs say little about a classifier, so I generated 150 PlainLang
programs from a fixed seed. They mix counted and uncounted loops, loop bounds hidden behind
function calls, counters aliased through other variables, waits inside conditionals, and
two-phase programs. Each was run \emph{alone} in the kernel, and its measured CPU time per
wait served as ground truth (Table~\ref{tab:acc}).

\begin{table}[t]
\centering\small
\begin{tabular}{@{}lrrr@{}}
\toprule
Analysis & Agreement & Too light & Too heavy \\ \midrule
basic & 62\% & 56 & 1 \\
counted & \textbf{76\%} & 36 & 0 \\
safe (heavy-biased) & 59\% & 11 & 50 \\ \bottomrule
\end{tabular}
\caption{Classifier vs.\ measured behavior, 150 programs. ``Too light'' means the guessed
class demands less CPU than measured.}\label{tab:acc}
\end{table}

The counted analysis was right 97\% of the time on counted \texttt{while} loops, and only
53--67\% on its documented blind spots. Just 2 of 150 were gross errors (CPU-bound
labeled interactive or vice versa). The finding I did not predict: \textbf{all 36 errors
went the same way}. Every wrong guess underestimated demand. Two rules cause this. Unknown
loops are assumed to run 100 times, and a wait inside a conditional is counted as if it
happened on every iteration.

Biasing the classifier toward heavy guesses (unknown loops assumed to run 10{,}000 times, a
conditional counting its fewest waits) did not fix this. It moved the errors to the other
direction and cut accuracy to 59\%. Three of its four pre-registered predictions failed. When
loop bounds span four orders of magnitude, no constant guess for an unknown bound is right.

\subsection{Who pays for a wrong hint?}
I ran ten mislabeled programs from the corpus next to a real interactive program and a CPU
hog (120 runs, predictions committed first). A \emph{too-heavy} hint cost only the
mislabeled program: under \pol{hint} its own p95 latency rose from about 1\us{} to exactly
100\ms{}, the aging threshold, while its neighbor stayed at 1\us{} in every case. A
\emph{too-light} hint delayed the neighbor by more than 1\ms{} in only 2 of 5 cases, and the
harm scaled with the program's burst length. Under \pol{hfair}, changing the hint changed
the neighbor's latency in none of the cases beyond the tick-collision floor. So the
classifier's characteristic error is the dangerous kind for hint-as-priority, and a harmless
kind for \pol{hfair}.

\section{Threats to validity}\label{sec:threats}
\begin{itemize}\setlength\itemsep{1pt}
\item \textbf{Emulated hardware.} Every instruction costs the same; there are no caches,
pipelines or TLBs. Only relative comparisons are meaningful.
\item \textbf{Synthetic workloads.} ``I/O'' is a timed sleep and ``user events'' arrive on
fixed periods. The classifier corpus is generated, not written by real users.
\item \textbf{An approximation of CFS.} \pol{fair} has no weights, nice levels, groups or
EEVDF deadlines~\cite{eevdf}, and runs on one CPU.
\item \textbf{The author wrote the baselines.} Constant sweeps (boost, aging, granularity,
sleeper credit, MLFQ depth, slice floor) showed some conclusions depend on such constants:
MLFQ's interactive latency ranged from 4.4 to 18.8\ms{} with its number of levels alone.
\item \textbf{Measurement apparatus.} A seed-for-seed regression check found that MLFQ's
boost timer carried state between runs, making MLFQ-family results depend on plan order (up
to 29\ms{} in one metric). It was fixed and the entire dataset rerun on one kernel image.
\item \textbf{Post-hoc elements.} The \pol{hfair} design and the tick-collision explanation
came after looking at data; both received separate pre-registered tests.
\end{itemize}

\section{Related work}
MLFQ originates with CTSS~\cite{corbato62} and is presented in standard texts~\cite{ostep}.
Proportional-share scheduling includes lottery and stride scheduling~\cite{lottery} and
virtual-time schedulers such as EEVDF~\cite{eevdf}, which Linux adopted after
CFS~\cite{cfs}. A systematic review of work on passing program-level hints to schedulers is
left to the author's ongoing literature review. This report makes no claim of novelty
until that review is complete.

\section{Conclusion}
Compiler hints look transformative against textbook schedulers, and that is the wrong
comparison. Against a scheduler that measures CPU usage directly, declared behavior adds
little to latency, because the scheduler already sees it. What a hint can add is
information about the \emph{future}: that a thread will want the CPU for a long time and
should not be switched out every millisecond. The robust way to use that information is to
let it change only an efficiency parameter and never priority. Then the classifier's
errors, which this study shows are systematic, cost efficiency rather than other programs'
responsiveness. Open next steps are real hardware, real user programs, and comparison with
the actual Linux scheduler.

\begin{thebibliography}{9}\small
\bibitem{corbato62} F.~J. Corbató, M.~M. Daggett, R.~C. Daley. An experimental time-sharing
system. \emph{AFIPS Spring Joint Computer Conference}, 1962.
\bibitem{ostep} R.~H. Arpaci-Dusseau, A.~C. Arpaci-Dusseau. \emph{Operating Systems: Three
Easy Pieces}. Arpaci-Dusseau Books.
\bibitem{lottery} C.~A. Waldspurger, W.~E. Weihl. Lottery scheduling: flexible
proportional-share resource management. \emph{OSDI}, 1994.
\bibitem{eevdf} I.~Stoica, H.~Abdel-Wahab. Earliest eligible virtual deadline first: a
flexible and accurate mechanism for proportional share resource allocation. Technical
report, Old Dominion University, 1995.
\bibitem{cfs} I.~Molnar. Completely Fair Scheduler. Linux kernel documentation
(\texttt{Documentation/scheduler/sched-design-CFS.rst}), 2007.
\end{thebibliography}

\end{document}
